\documentclass{article}

\usepackage[preprint]{neurips_2024}

\usepackage[utf8]{inputenc}
\usepackage[T1]{fontenc}
\usepackage{hyperref}
\usepackage{url}
\usepackage{booktabs}
\usepackage{amsfonts}
\usepackage{nicefrac}
\usepackage{microtype}
\usepackage{xcolor}
\usepackage{graphicx}
\usepackage{listings}
\usepackage{amsmath}

\lstdefinestyle{yamlstyle}{
  basicstyle=\ttfamily\small,
  columns=fullflexible,
  breaklines=true,
  frame=single,
  framesep=3pt,
}

\title{Deliverable Contracts: A Deterministic Governance Layer for\\
Long-Horizon Agent Tasks}

\author{
  Anonymous Author(s) \\
  Long-Horizon Agent Governance Research Project
}

\begin{document}

\maketitle

\begin{abstract}
Long-horizon autonomous agent runs are typically governed by trusting the
agent's own self-report that a task is ``done.'' We show, on 300 real,
publicly published coding-agent submissions to SWE-bench Lite, that this
status-quo governance signal has zero ability to distinguish a sound
deliverable from an unsound one: every instance the agent's own harness
reported as ``submitted'' is treated identically by this baseline,
regardless of whether the underlying patch is structurally coherent. We
introduce \textsc{gatekeep}, a deterministic, offline, sub-millisecond
\emph{deliverable contract} checker: a small declarative schema for what a
long-horizon task must produce, plus a checker that verifies the actual
artifact state against that schema without any model call. Applied to the
same 300 real agent submissions and scored against the official
SWE-bench evaluation harness's ground truth, \textsc{gatekeep} achieves
perfect recall (1.00) and improves precision over the status-quo baseline
(0.2438 vs.\ 0.2404) by flagging four deliverables as structurally
unsound before any test suite is run, at zero added false negatives.
We report the full precision/recall/F1 table, document two real
false-positive bugs we found and fixed during the experiment, and discuss
honestly what a cheap deterministic check can and cannot catch. The system,
benchmark scoring code, and raw logs are publicly released.
\end{abstract}

\section{Introduction}

Long-horizon agent tasks --- multi-hour autonomous coding sessions,
research tasks, or any workflow where an LLM-driven agent operates for
extended periods with multiple intermediate steps before producing a final
set of deliverables --- have a specific, under-discussed failure mode that
we call \emph{deliverable drift}: the task specifies $N$ required outputs,
and over the course of a long run the agent silently drops, truncates, or
fabricates one or more of them while still reporting overall success.

The standard way agent harnesses today decide whether to trust a
deliverable is to trust the agent's own completion signal: an exit code,
an ``I'm done'' message, or a non-null artifact in the expected location.
We call this \emph{status-quo governance}. It is cheap and ubiquitous, but
it has an obvious weakness: it asks the agent to grade its own homework.

This paper makes two contributions:

\begin{enumerate}
  \item \textbf{A measurement.} Using 300 real, publicly available agent
  submissions to SWE-bench Lite (SWE-agent driving Claude 3.5 Sonnet,
  published by the SWE-bench project), we measure how often status-quo
  governance --- ``a non-empty submission with a clean exit status'' ---
  is wrong about deliverable soundness, against real ground truth from the
  official SWE-bench evaluation harness. We find it has \emph{zero}
  discriminative power on this population: every instance with a non-null,
  cleanly-submitted patch is classified identically as ``shippable,''
  whether or not the patch actually resolves the issue (Table~\ref{tab:ab}).

  \item \textbf{A system.} \textsc{gatekeep}, a deterministic deliverable
  governance layer. Long-horizon task owners declare a \emph{deliverable
  contract} (a YAML file: what must exist, in what shape, with what
  content properties) and run a sub-millisecond, offline checker against
  it. We show this catches a measurable subset of unsound deliverables
  that status-quo governance misses, with no added false negatives, using
  only deterministic, reproducible checks --- no LLM call, no API key, no
  network dependency.
\end{enumerate}

We are explicit about scope: \textsc{gatekeep} is not a smarter SWE-bench
solver, and a deterministic syntactic/structural check cannot, by
construction, verify semantic correctness (whether the patch actually
fixes the bug). What it verifies is a necessary but not sufficient
condition: did the agent produce something that even has the right shape
to be correct? Section~\ref{sec:limitations} discusses this honestly,
including two real false-positive bugs we discovered and fixed while
running the experiment in this paper.

\section{Related Work}

\textbf{Agent benchmarks.} SWE-bench~\citep{jimenez2024swebench} evaluates
whether an LM-driven agent can resolve real GitHub issues, scored by
running the project's test suite against the agent's patch in a sandboxed
container. SWE-bench Lite is a 300-instance subset used for faster
iteration. The SWE-bench project also maintains a public
\texttt{experiments} repository of raw submissions (trajectories, patches,
and evaluation reports) from many agents evaluated against the benchmark;
we use one such public, externally-produced submission
(SWE-agent~\citep{yang2024sweagent} driving Claude 3.5 Sonnet) as the
substrate for our measurement, rather than producing new rollouts
ourselves.

\textbf{LLM-as-judge evaluation.} A large body of agent evaluation work
uses an LLM to grade another LLM's output against a rubric. This is
expressive but non-deterministic and requires a model call (cost, latency,
API access) per judgment. \textsc{gatekeep} takes the opposite design
point: every check is a pure function over the filesystem, so the same
contract gives the same verdict for anyone, forever, independent of model
version or API availability. We do not claim this is strictly better for
every use case --- only that it is a different, complementary point in the
design space, useful precisely when reproducibility and zero marginal cost
matter more than rubric expressiveness.

\textbf{CI/build verification.} Software engineering has long used
deterministic gates (linters, type checkers, test runners) to verify build
artifacts before merge. \textsc{gatekeep} applies the same philosophy to
agent task \emph{deliverables} broadly (not just code): a declarative
contract plus a fast, auditable checker, designed to be dropped into
existing CI pipelines (we ship a GitHub Actions wrapper).

\section{The gatekeep System}

\subsection{Deliverable contracts}

A contract (Listing~\ref{lst:contract}) declares a list of named
\emph{deliverables}. Each deliverable has a severity (\texttt{required} or
\texttt{advisory}) and a list of \emph{checks}. A run passes iff every
\texttt{required} deliverable's checks all pass; \texttt{advisory}
failures are reported but never block.

\begin{lstlisting}[style=yamlstyle, caption={A deliverable contract.}, label={lst:contract}]
name: my-project-deliverables
deliverables:
  - id: api-patch
    description: "the agent's code change is real"
    severity: required
    checks:
      - kind: valid_unified_diff
        path: "patch.diff"
        forbid_test_only: true
  - id: readme
    severity: required
    checks:
      - kind: min_words
        path: "README.md"
        min_words: 50
      - kind: no_placeholder_text
        path: "README.md"
\end{lstlisting}

\subsection{Check kinds}

The current engine implements eight deterministic check kinds:
\texttt{file\_exists} (glob match with a minimum count),
\texttt{min\_size} (byte-size floor), \texttt{min\_words} (word-count
floor), \texttt{no\_placeholder\_text} (regex bank for TODO/FIXME/lorem
ipsum/``not implemented''/etc., with an \texttt{allow}-list escape hatch
and an optional diff-aware mode that scans only \emph{added} lines),
\texttt{json\_valid}, \texttt{regex\_required}, \texttt{valid\_unified\_diff}
(structural diff validation: real \texttt{@@} hunks, minimum file count,
optional ``no test-only patches'' rule), and \texttt{forbid\_glob} (ban
stray scratch files). Every check is a small, independently testable pure
function; adding a new kind is one function plus a registry entry.

\subsection{Two install paths, kept in parity}

\textsc{gatekeep} ships as (a) a pip-installable package with a
\texttt{gatekeep} console entry point, and (b) a single, zero-dependency
file (\texttt{gatekeep\_single.py}, stdlib only, including a self-contained
YAML-subset parser) that a stranger can \texttt{curl} and run with no
install step. Both paths read the identical contract format and are kept
in lock-step by an automated parity test suite
(\texttt{test\_single\_file\_parity.py}) that runs the same contracts
through both implementations and asserts identical verdicts. The full test
suite (22 test executions at submission time) covers both engines plus the YAML
subset parser, including regression tests for every bug found during this
paper's experiment (Section~\ref{sec:limitations}).

\subsection{Self-application}

As a qualitative check, we wrote a \textsc{gatekeep} contract describing
this project's own five required deliverables (deployable system, paper,
public-benchmark A/B, marketing plan, landing page) and ran
\textsc{gatekeep} against this very repository before submission. This
dogfooding step is what surfaced two of the bugs fixed in
Section~\ref{sec:limitations}: the contract's own README initially failed
its own placeholder-text check because the README \emph{documents} the
words ``TODO''/``FIXME'' as part of describing the feature --- a clean
self-referential edge case, resolved via the contract's \texttt{allow}
escape hatch rather than by weakening the check.

\section{Experimental Setup}

\subsection{Substrate}

We use SWE-bench Lite~\citep{jimenez2024swebench}, a 300-instance subset
of SWE-bench covering real GitHub issues across 11 popular Python
repositories. Rather than producing new agent rollouts ourselves, we use a
real, publicly published, externally produced submission from the
official SWE-bench \texttt{experiments} repository
(\url{https://github.com/swe-bench/experiments}): run
\texttt{20240620\_sweagent\_claude3.5sonnet}, SWE-agent~\citep{yang2024sweagent}
driving Claude 3.5 Sonnet. For every instance this run publishes (in a
public, unauthenticated Amazon S3 bucket) the agent's full step-by-step
trajectory, its final submitted unified diff (\texttt{patch.diff}), and a
ground-truth evaluation report (\texttt{report.json}) produced by
princeton-nlp's own SWE-bench evaluation harness, independent of this
work. We fetched all three artifacts for all 300 instances; 287 have a
usable ground-truth report (the remaining 13 either never produced a
generation at all, 11 instances, or are missing a report for unrelated
reasons, 2 instances --- see \texttt{benchmark/data/fetch\_summary.json}).

\subsection{Arms}

\textbf{Arm A (status-quo baseline).} An instance is ``shippable'' iff the
agent's trajectory contains a non-empty \texttt{submission} string and the
trajectory's \texttt{exit\_status} begins with \texttt{"submitted"}. This
mirrors how most agent harnesses gate today: trust the agent's own
completion signal, with no inspection of the artifact's content.

\textbf{Arm B (gatekeep).} An instance is ``shippable'' iff
\textsc{gatekeep}'s deterministic contract
(\texttt{benchmark/contracts/swebench\_patch.yml}) passes against that
instance's \texttt{patch.diff}: the patch must exist, be at least 30
bytes, be a syntactically valid unified diff touching at least one
non-test file, and contain no placeholder/error markers
(\texttt{NotImplementedError}, \texttt{FIXME}, etc.) on any line the patch
\emph{adds}. Both arms make their decision before consulting
\texttt{report.json}.

\subsection{Scoring}

We treat ``shippable'' as a binary classifier for the real ground-truth
label \texttt{resolved} from \texttt{report.json}, and report precision,
recall, F1, and accuracy for both arms over the 287 instances with ground
truth, plus a headline operational number: of the 218 instances that were
\emph{not} actually resolved (``doomed'' deliverables), how many did each
arm let through to the costly test-execution stage versus catch for free
beforehand.

\section{Results}
\label{sec:results}

Table~\ref{tab:ab} reports the full confusion-matrix metrics for both arms
on the 287 SWE-bench Lite instances with ground truth, computed by
\texttt{benchmark/scripts/run\_ab.py} and saved verbatim to
\texttt{benchmark/results/ab\_summary.json} and
\texttt{ab\_raw\_rows.json} (one row per instance, with the exact check
output that produced each verdict).

\begin{table}[t]
  \caption{Arm A (status-quo self-report governance) vs. Arm B
  (\textsc{gatekeep}) on 287 real SWE-bench Lite instances, scored against
  official SWE-bench ground truth. TP/FP/FN/TN are with respect to
  predicting \texttt{resolved=True}.}
  \label{tab:ab}
  \centering
  \begin{tabular}{lrrrrrrrr}
    \toprule
    Arm & TP & FP & FN & TN & Precision & Recall & F1 & Accuracy \\
    \midrule
    A: status-quo baseline & 69 & 218 & 0 & 0 & 0.2404 & 1.0000 & 0.3876 & 0.2404 \\
    B: gatekeep            & 69 & 214 & 0 & 4 & 0.2438 & 1.0000 & 0.3920 & 0.2544 \\
    \bottomrule
  \end{tabular}
\end{table}

Out of 287 instances, 69 were genuinely resolved and 218 were not.
\textbf{Arm A has zero discriminative power}: TN $=0$, meaning that among
\emph{every} instance the agent's own harness considered cleanly submitted,
not a single one was flagged as suspect, regardless of whether the
underlying patch was sound. This is the quantitative version of ``the
agent always says it's done.''

\textbf{Arm B catches 4 of the 218 doomed deliverables for free}, before
any test suite runs, with \textbf{zero additional false negatives}
(recall stays at 1.00 -- it never incorrectly flags a deliverable that
was actually resolved). The four catches break down as: three instances
where the agent's submitted patch touched \emph{only} test files (e.g.\
\texttt{django\_\_django-16229}, whose patch modifies
\texttt{tests/custom\_test\_runner.py}, \texttt{tests/test\_arrayfield\_\allowbreak admin.py}, and two other test files but no source file at all
--- a patch that, by construction, cannot fix the underlying issue), and
one instance (\texttt{sympy\_\_sympy-22005}) where the submitted diff adds
a \texttt{raise NotImplementedError(...)} for an unsupported code path.

The deterministic check itself is fast: scoring all 300 local instances
(running the full contract against each instance's artifact directory)
takes 0.157\,s wall-clock total, or 0.52\,ms per instance, on a single
consumer laptop core, with no network call.

\subsection{Where the four catches come from, qualitatively}

We manually inspected the four instances Arm B flags and Arm A does not
(Table~\ref{tab:catches}). All four are genuine: none is an artifact of an
overly aggressive check rule (we verified this by also inspecting, and
ruling out as false positives, two close calls discussed in
Section~\ref{sec:limitations}).

\begin{table}[h]
  \caption{The four doomed deliverables \textsc{gatekeep} flags that the
  status-quo baseline does not.}
  \label{tab:catches}
  \centering
  \small
  \begin{tabular}{lp{8.7cm}}
    \toprule
    Instance & Reason flagged \\
    \midrule
    \texttt{django\_\_django-16229} & Patch touches only test files (4 files, 0 source files). \\
    \texttt{django\_\_django-16408} & Patch touches only \texttt{tests/known\_related\_objects/tests.py}. \\
    \texttt{pytest-dev\_\_pytest-7168} & Patch touches only \texttt{test\_repr\_error.py}. \\
    \texttt{sympy\_\_sympy-22005} & Patch adds \texttt{raise NotImplementedError(...)} for an unhandled case. \\
    \bottomrule
  \end{tabular}
\end{table}

\section{Limitations and Honest Failure Modes}
\label{sec:limitations}

We report two real bugs found and fixed during this experiment, because a
governance tool's credibility depends on being honest about its own false
positives.

\textbf{Bug 1: substring test-path matching.} An early version of the
``forbid test-only patches'' rule classified a path as a test file if the
substring \texttt{"test"} appeared anywhere in the path. This
misclassified \texttt{src/\_pytest/assertion/rewrite.py} (a real source
file in the \texttt{pytest} project) as a test file purely because
\texttt{"pytest"} contains the substring \texttt{"test"}, producing a
false negative on instance \texttt{pytest-dev\_\_pytest-11143} (a resolved
instance that Arm B would have incorrectly flagged as unsound). We fixed
this with a path-component-aware predicate (a component must equal
\texttt{test}/\texttt{tests}, or the filename stem must start with
\texttt{test\_} or end with \texttt{\_test}) and added a regression test.

\textbf{Bug 2: removed-line placeholder false positives.} The
placeholder-text scanner initially scanned the entire diff text, including
context and \emph{removed} lines. This flagged
\texttt{django\_\_django-12286} (a resolved instance) because its patch
\emph{deletes} two pre-existing \texttt{\# FIXME} comments as part of
removing dead code --- the agent did not introduce a placeholder, it
removed two. We added a \texttt{diff\_added\_lines\_only} mode that scans
only lines the patch adds, and a regression test asserting a removed
\texttt{FIXME} does not trigger a failure while an added one still does.

Both bugs were caught by running the real experiment end-to-end against
real data and inspecting disagreements between gatekeep and ground truth
--- not by code review alone. We consider this evidence in favor of the
project's own thesis: deterministic checks are only as trustworthy as the
checks you write, and validating them against real outcomes (not just unit
tests of the checker's own logic) is necessary.

\textbf{What deterministic checks cannot do.} Even after both fixes, Arm
B's precision (0.244) is far below 1.0: 214 of its 283 ``shippable'' calls
are deliverables that are not actually resolved. This is expected and by
design --- a syntactically valid, non-placeholder, source-touching patch
can still be semantically wrong. \textsc{gatekeep} verifies a necessary,
not sufficient, condition for deliverable soundness. It is not a
replacement for running the actual test suite or for human review; it is
a cheap, instant pre-filter that catches the cases where the deliverable
does not even have the right \emph{shape} to be correct, before paying the
cost of deeper verification.

\textbf{One more honest caveat.} Our \texttt{NotImplementedError} pattern
flags any added line containing that token, including legitimate
exception-handling code (e.g.\ \texttt{except NotImplementedError:},
\texttt{assertRaises(NotImplementedError)}) alongside genuine unfinished
stubs. In our data this did not produce a false positive on a resolved
instance (the one match, \texttt{sympy\_\_sympy-22005}, is a true positive
for ``doomed'' regardless of the specific line that triggered it), but a
contract author relying on this pattern in a codebase with heavy
exception-handling idioms should expect lower precision on that specific
check and may want to narrow the pattern (e.g.\ to
\texttt{raise NotImplementedError} only, excluding
\texttt{except}/\texttt{assertRaises} contexts).

\section{Discussion}

The headline finding of this paper is not that \textsc{gatekeep} is a
large improvement over the baseline in aggregate F1 (0.392 vs.\ 0.388 ---
a modest difference, because most of the discriminative work on this
particular dataset would require semantic, not structural, verification).
The headline finding is that \textbf{the status-quo baseline has
\emph{zero} true negatives on real data}: it provides no information
beyond ``did the agent's own harness claim success,'' which is exactly the
self-report problem this paper set out to quantify. Any non-zero,
zero-false-negative improvement over a classifier with literally no
discriminative power is operationally meaningful, because it is pure
signal added at effectively zero cost (0.52\,ms/instance, no model call,
no network, fully reproducible).

We see \textsc{gatekeep} as one layer in a larger governance stack, not a
complete solution: deterministic structural checks for the necessary
conditions (this work), optionally composed with semantic checks (test
execution, or an LLM-judge plugin for genuinely soft rubric criteria) for
the sufficient conditions. The contribution of this paper is showing that
the cheap layer is not vacuous --- it catches real, qualitatively
inspectable failures on real public data --- and that it can be built,
tested, and shipped as a small, auditable, dependency-light artifact that
a long-horizon agent task owner can adopt in minutes.

\section{Reproducibility}

All numbers in this paper are produced by code in the public repository
accompanying this submission. \texttt{benchmark/scripts/fetch\_swebench\_run.py}
re-downloads the 300 instances' raw artifacts from the public,
unauthenticated S3 bucket; \texttt{benchmark/scripts/distill\_trajs.py}
extracts the trajectory summary used for Arm A;
\texttt{benchmark/scripts/run\_ab.py} reproduces Table~\ref{tab:ab} and
both raw-row and summary JSON files exactly. The \textsc{gatekeep} engine
itself has 22 unit/regression/parity test executions
(\texttt{gatekeep/tests/}), all passing at submission time, including a
dedicated regression test for each bug in Section~\ref{sec:limitations}.
No number in this paper was generated by, or required, an LLM API call.

\bibliographystyle{plainnat}
\begin{thebibliography}{9}

\bibitem[Jimenez et al., 2024]{jimenez2024swebench}
Carlos E. Jimenez, John Yang, Alexander Wettig, Shunyu Yao, Kexin Pei,
Ofir Press, and Karthik Narasimhan.
\textit{SWE-bench: Can Language Models Resolve Real-World GitHub Issues?}
International Conference on Learning Representations (ICLR), 2024.

\bibitem[Yang et al., 2024]{yang2024sweagent}
John Yang, Carlos E. Jimenez, Alexander Wettig, Kilian Lieret, Shunyu Yao,
Karthik Narasimhan, and Ofir Press.
\textit{SWE-agent: Agent-Computer Interfaces Enable Automated Software
Engineering.}
arXiv:2405.15793, 2024.

\end{thebibliography}

\end{document}
